% supplement.tex -- supplementary material for the TNNLS submission (uploaded as a separate PDF labelled "supplementary").
% Excluded from the 15-page limit (CIS Information for Authors, section H). Accepted formats: TXT, DOC(X), PDF (Author Center).
% Reviewers are not required to read it: the main paper must stand alone. Anonymous: no author names or links.
% Build: see ../build.sh  (latexmk -pdf -cd supplement/supplement.tex)
\documentclass[journal]{IEEEtran}
\usepackage[T1]{fontenc}
\usepackage{lmodern}
\usepackage{amsmath,amssymb,amsthm}
\usepackage{graphicx}
\usepackage{booktabs}
\usepackage{xcolor}
\usepackage[ruled,vlined,linesnumbered]{algorithm2e}
\usepackage{microtype}
\usepackage[colorlinks=true,linkcolor=blue,citecolor=blue,urlcolor=blue,breaklinks=true]{hyperref}
\usepackage[capitalise]{cleveref}
\newcommand{\todo}[1]{{\color{red}\textbf{[#1]}}}
\renewcommand{\thesection}{S\arabic{section}}   % S1, S2, ... so references to the main paper stay unambiguous
\renewcommand{\thetable}{S\arabic{table}}
\renewcommand{\thefigure}{S\arabic{figure}}
\renewcommand{\theequation}{S\arabic{equation}}

\begin{document}
\title{Supplementary Material for [AUTHOR: title]}
\author{Anonymous authors}
\maketitle

\section{Additional Results}
% TODO(supplement): comparison tables, ablations, per-corruption breakdowns, additional algorithm boxes, extended proofs,
% protocol and deviation log. Cross-reference the main paper as "Section IV of the main paper" in words (no \cref across files).
\todo{TODO: supplementary content}

\section{Reproducibility}
% TODO(supplement): environment, seeds, commands, README content (description, size, platform, environment, components, setup,
% run instructions, output description, contact -- Author Center "Creating a README File"). Anonymised links only.
\todo{TODO: reproducibility details}

\end{document}
